\section{Docker}
Docker ist aus der Welt der Container nicht wegzudenken. 
Zwar hat diese Firma die Container nicht erfunden, aber sie so weit verbreitet, dass sie als Synonym dafür benutzt werden, wie Lego für Klemmbrettbausteine

Ein Container wird durch ein Image erstellt. Die Besonderheit dabei ist, dass das Image ein Read-only-Dateisystem ist und daher im laufenden Container nie verändert werden kann.
Um ein Image zu erstellen, muss in Docker zunächst ein Dockerfile angelegt werden.

Eine weitere Besonderheit bei Containern sind die sogenannten Volumes. Bei Containern werden Dateien direkt im Host-System gespeichert. Diese Verzeichnisse sind vom Container getrennt.
Dadurch bleiben die Daten auch erhalten, wenn der Container neu eingerichtet wird.
Docker legt dabei normalerweise automatisch ein Verzeichnis an, wenn ein Container ein Volume benötigt.

Images müssen aber nicht immer, selber erstellt werden, sondern können auch extern aus dem Internet gezogen werden.
Dabei bietet sich bei Docker besonderes der Docker Hub an. In dieser von Docker betriebenen Platform können viele Image gefunden werden die von der Docker-Community erstellt wurden.
Dies kann wie ein Art Werkzeugkoffer angesehen werden, von dem vorgefertige Anwendungen und Umgebungen genutzt werden können oder sogar mit dem eigenen Image kombinierbar sind,
wie z.B. mongodb, nginx und so weiter \parencite[vgl. 50]{kofler2023docker}. 

Natürlich gibt es da auch wieder mehrere alternativen für den Docker Hub.